% !Mode:: "TeX:UTF-8"
\fancyhead[CO]{\song\wuhao ~\chaptername \quad 相关知识及理论介绍~}
\fancypagestyle{plain}{\fancyhead[CO]{\song\wuhao ~\chaptername \quad 相关知识及理论介绍~}}

\chapter{相关知识及理论介绍}

% 第一节
\section{异构计算加速平台}
\subsection{异构计算概述}

随着深度学习技术快速发展，尤其是 Transformer 大语言模型参数量达到万亿规模，传统以 CPU 作为核心的同构计算模式已经难以满足 AI 任务在计算和访存上的需求，面向人工智能的异构计算因此成为支撑现代深度学习模型训练和推理的基础设施。深度神经网络和 Transformer 模型表现出高度规则化及数据并行化的特征，计算架构由此从以控制为中心的通用处理器 CPU，转变为以数据为中心的异构协同计算模式。CPU 在该模式中主要负责逻辑控制、数据预处理和任务调度，加速器则专门进行矩阵乘法、卷积、注意力等主要算子的执行\cite{jouppi2017datacenter}。

为争取大模型时代的算力主导权，全球半导体厂商推出多种技术路线的异构加速平台，形成多样但生态割裂的局面。主流平台主要包括通用路线和专用人工智能芯片路线\cite{jouppi2021ten}。具有 CUDA 生态的 NVIDIA GPU 作为通用 GPU，已经成为当下主流且成熟的硬件载体。AMD GPU 和国产海光 DCU 也使用相似的架构，它们是形成国产自主算力基础的重要力量。专用人工智能芯片路线包括华为昇腾 NPU、Google TPU 和寒武纪 MLU 等，专门对深度学习计算进行定制，以获得更好的能效比。不过，各类平台在指令集架构、微架构设计、存储层级和编程模型上存在明显差异，上层深度学习框架因此很难通过统一后端充分利用不同硬件的性能。在各种方案中，NVIDIA GPU 代表工业界性能基准和优化标准，海光 DCU 体现国产替代的主要方向和典型架构。分析算子性能瓶颈、设计跨平台编译优化策略，都需要以理解这两类代表性硬件的架构及编程模型为前提。

综上所述，异构计算平台为 Transformer 关键算子的高性能执行提供了重要硬件基础，但不同平台之间的执行模型、线程组织方式和存储层次差异，也给算子跨平台生成和优化带来了直接挑战。因此，本文后续围绕 NVIDIA GPU 和海光 DCU 两类具有代表性的异构平台展开分析，为后续编译优化方法的设计提供硬件背景。

\subsection{NVIDIA GPU 架构及编程模型}

随着深度学习领域不断发展，NVIDIA GPU 凭借良好的并行计算能力和完善的 CUDA 软件生态，已经成为模型训练和推理的通用平台。其硬件架构从 Volta\cite{nvidia_volta} 和 Ampere 发展到 Hopper、Blackwell 等架构，并且持续加入适用于 AI 负载的专用优化特性。图 \ref{2_1} 展示了配备 144 个流多处理器（Streaming Multiprocessor, SM）的完整 NVIDIA H100 GPU 核心。

\begin{figure}[htbp]
\centering
\includegraphics[width=0.95\textwidth]{2_1}
\caption{NVIDIA H100 GPU 核心\cite{nvidia_hopper}}\label{2_1}
\vspace{\baselineskip}
\end{figure}

硬件架构方面，GPU 的主要计算单元为 SM，图 \ref{2_2} 展示了 NVIDIA H100 GPU 的 SM 架构。该 SM 内部包括大量 CUDA 核心用于进行标量运算，同时也设有为深度学习定制的张量核心。张量核心能够高效执行混合精度的矩阵乘累加运算，是 Transformer 模型中多头注意力机制和矩阵乘法的核心加速引擎。

GPU 采用显式的多级存储机制\cite{mei2016dissecting}来解决内存瓶颈问题，该机制包括速度较快但容量较少的寄存器堆、可当作可编程暂存空间使用的共享内存和具有较大容量但延迟较高的全局显存。高效利用该存储层次，特别是通过分块技术来优化共享内存的访问，是达到高性能计算的主要方式\cite{nickolls2008scalable}。

在编程模型层面，CUDA 采用单指令多线程模型，通过大规模线程并行来掩盖访存延迟。如图 \ref{2_3} 所示，在 NVIDIA GPU 编程模型中，其线程组织呈现为“网格-线程块-线程”的三级逻辑抽象，其中线程簇（Cluster）层级为 Hopper 架构后特有抽象层级，不属于传统编程模型的层次结构。

处于同一个线程簇内的线程块可以进行数据通信，这样能够减少对 HBM 的访问。线程束（Warp）由 32 个线程组成，是硬件执行的基本单位。同一个线程束内部的线程利用锁步机制运行相同指令，条件分支的出现会引起分支发散，对执行效率造成严重影响。线程块作为资源分配和调度的主要单位，会被映射到 SM 上执行。SM 中的线程束调度器利用动态上下文切换方式，在等待数据的线程束和准备计算的线程束之间进行快速切换，实现计算与访存重叠，进而掩盖延迟。因而，保持足够的占用率，也就是在 SM 上驻留足够多的活跃线程束，从而使计算流水线维持饱和状态，是让 GPU 达到最佳性能的重要原则\cite{baghsorkhi2010adaptive}。

\begin{figure}[htbp]
\centering
\includegraphics[width=0.75\textwidth]{2_2}
\caption{NVIDIA H100 GPU SM 架构\cite{nvidia_hopper}}\label{2_2}
\vspace{\baselineskip}
\end{figure}

\begin{figure}[htbp]
\centering
\includegraphics[width=0.55\textwidth]{2_3}
\caption{NVIDIA GPU 编程模型的逻辑抽象\cite{nvidia_cuda_programming_guide}}\label{2_3}
\vspace{\baselineskip}
\end{figure}

算子开始运行之后，全局调度器会把网格（Grid）里的线程块分配给各个 SM。在 SM 内部，线程束调度器从处于就绪状态的线程束中选择指令进行发射。该设计利用上下文切换机制来达到延迟隐藏的目的，当某个线程束在等待全局显存数据时，调度器会立刻切换到另一个数据已经准备好的线程束进行计算。因此，为了充分利用 GPU 的计算资源，必须保持较高的占用率，也就是在 SM 上驻留足够多的活跃线程束，使计算流水线持续运行。充分了解 NVIDIA GPU 在 SM 微架构、多级存储、SIMT 执行模型上的软硬件协同设计原理，构成了对底层算子性能进行分析和优化的必要基础。

\subsection{海光 DCU 架构及编程模型}

海光 DCU 是国产异构计算领域中有代表性的硬件平台，通常用于高性能计算中心和大规模人工智能算力集群。其设计理念继承了通用 GPU 的大规模并行架构，并且采用 SIMT 执行模型，在微架构实现和资源组织方式上与 NVIDIA GPU 存在差别，该差异成为跨平台算子优化需要考量的主要因素。
\begin{figure}[htbp]
\centering
\includegraphics[width=0.95\textwidth]{2_4}
\caption{DCU Z100 (CNDA1.0) CU 架构\cite{amd_cdna}}\label{2_4}
\vspace{\baselineskip}
\end{figure}

海光 DCU 在硬件架构上得到了 AMD CDNA 架构的技术授权，所以它的架构特点和 AMD CDNA 架构\cite{amd_cdna}大体一致，其核心计算模块是计算单元（Compute Unit, CU），在功能上同 NVIDIA 的 SM 相对应。图 \ref{2_4} 展示了 CU 的架构。CU 内部的标量单元、矢量单元、专用矩阵乘法加速单元（Matrix Core）被集成在一起，以此支持 AI 混合精度计算。DCU 的基本调度单位是包括 64 个线程的波前（Wavefront），这与 NVIDIA GPU 不同，NVIDIA 的线程束有 32 个线程。若继续使用 32 线程粒度，直接把算法移植到 DCU 会引起线程空闲或者同步错误，进而让硬件利用率下降。

海光 DCU 在存储架构上使用了多级存储层次结构，并且把片上共享内存称为 LDS（Local Data Share）。LDS 通常分为 32 个 Bank，每个 Bank 的宽度为 4 字节。因为 DCU 的 Wavefront 宽度是 64，而 LDS 的 Bank 数量只有 32，在一次并发访问中，发生 Bank 冲突的概率和严重程度比 NVIDIA 平台更高。例如，在步长为 32 的跨步访问模式下，NVIDIA 平台不会产生冲突，但是 DCU 会出现串行化延迟。Access Scatter 优化策略的主要动机正是这一硬件特性。

在编程模型上，为支持 DCU 的应用生态，海光发布了全栈式异构计算开发包（DCU Toolkit, DTK）。它对标 NVIDIA CUDA Toolkit，并根据 AMD ROCm 开源生态进行定制优化，尝试作为连接上层应用和底层硬件的统一桥梁。就编程接口而言，DTK 把 HIP 当作主要编程语言，它的语法和 CUDA 非常相似，并且可以对代码进行 hipify 转换，而底层执行语义则完全适用于 DCU 的 Wavefront64 架构。在计算库支持上，DTK 集成了 rocBLAS 和 MIOpen 等官方加速库，这些库都经过了手工汇编优化。但是，由于算子融合能力有所不足并且版本迭代存在滞后，本文通过编译优化方式尝试超越该基准。在指令集架构方面，DCU 对 AMD 的 CDNA 指令集进行了扩展，部分设备也自行定义了硬件标识符，LLVM\cite{lattner2004llvm} 生态很难直接找到对应的指令集架构来完成编译工作，这直接增加了底层编译和优化的复杂程度。

% 第二节
\section{Transformer 模型架构与计算特性}

\subsection{Transformer 架构综述}
Transformer 模型最早由 Vaswani 等人在 2017 年的经典论文《Attention Is All You Need》\cite{vaswani2017attention}中提出。与传统的循环神经网络（RNN）和卷积神经网络（CNN）不同，Transformer 彻底抛弃了循环与卷积结构，完全基于注意力机制来捕捉序列中长距离的依赖关系。这种设计不仅解决了长序列训练\cite{press2021train}中的梯度消失问题，更极大地提升了模型的并行计算能力，使其能够充分利用现代 GPU 等异构硬件的计算资源。

\begin{figure}[htbp]
\centering
\includegraphics[width=0.5\textwidth]{2_5}
\caption{Transformer 架构\cite{vaswani2017attention}}\label{2_5}
\vspace{\baselineskip}
\end{figure}

图 \ref{2_5} 展示了采用编码器-解码器（Encoder-Decoder）模式\cite{xiong2020layer}的标准 Transformer 架构，它由若干结构相同的子层连接组成，一般记作 N 层堆叠结构。每一层内部利用残差连接\cite{he2016deep}（Residual Connection）和层归一化\cite{ba2016layer}（Layer Normalization）来进行稳定训练，使得在深层网络中可以较好地缓解梯度消失和数值不稳定问题。

Transformer 需要显式地引入位置信息，通过把输入符号嵌入向量同位置编码进行逐元素相加，来得到最终的输入表示\cite{su2024roformer}。编码器每层包括多头注意力机制和前馈全连接网络\cite{bebis2002feed}（Feed-Forward Network, FFN）。多头注意力机制将输入映射为查询、键、值向量并计算相关性，并行处理多个注意力头，然后拼接输出并进行线性变换。前馈全连接网络包括两层线性变换和非线性激活函数，是模型里计算量非常大的模块。解码器和编码器结构比较类似，并且加入编码器-解码器注意力模块，还通过掩码（Mask）多头自注意力来防止自回归生成时的信息泄漏。计算上，模型以矩阵乘法和注意力机制为主，注意力机制在长序列下具有 $O(L^{2})$ 复杂度，对硬件存储要求高，各模块的数据依赖特性为优化提供空间，也对异构硬件加速和编译器优化提出了更高要求。

\subsection{模型关键算子介绍}

在 Transformer 模型中，大部分计算负载集中在矩阵乘法和注意力机制这两个主要算子上。这两部分构成了模型的基本运算单元，它们在计算量和访存开销方面也占据主要位置，是影响模型推理效果的重要因素。图 \ref{2_6} 展示了矩阵乘法和注意力机制运算过程的示意图。因此，对这两个算子的高性能实现进行研究具有重要价值。
\begin{figure}[htbp]
\centering
\includegraphics[width=0.9\textwidth]{2_6}
\caption{矩阵乘法（左）和注意力机制（右）计算流程示意图}\label{2_6}
\vspace{\baselineskip}
\end{figure}


首先，矩阵乘法是 Transformer 中最基础且最频繁调用的线性代数操作。其计算形式可表示为：
\begin{equation}\label{matmul}
C_{ij} =\sum_{k=0}^{K-1} A_{ik}B_{kj}
\end{equation}
其中，$A\in \mathbb{R}^{M\times K}$、$B\in \mathbb{R}^{K\times N}$，输出矩阵 $C\in \mathbb{R}^{M\times N}$。
公式表明矩阵乘法是典型的计算密集型算子。在 Transformer 模型中，矩阵乘法广泛存在于注意力机制内部、前馈神经网络中的线性变换层以及全连接操作中。大模型通常具有较高的隐藏维度和批大小，所以矩阵乘法在整个推理阶段占据主要开销。矩阵乘法的性能高度依赖于分块策略、缓存复用及并行映射方式，因此成为深度学习编译优化的核心对象。

其次，注意力机制是 Transformer 模型的核心创新，其功能是为序列中每个位置动态建模与其他位置之间的依赖关系。其中缩放点积注意力（Scaled Dot-Product Attention）计算流程可表示为：
\begin{equation}\label{attention}
Attention \left ( Q,K,V \right ) = Softmax \left ( \frac{QK^{T} }{\sqrt{d_{k} } }  \right ) V
\end{equation}
其中 $Q$、$K$ 和 $V$ 分别为查询、键和值，$d_{k}$ 为键向量维度。从计算流程上看，多头注意力机制可分为三个阶段，首先计算相关性矩阵 $S=QK^{\mathrm{T}}$，然后对每一行执行缩放与 Softmax 归一化，最后将归一化结果与 $V$ 相乘得到输出结果。
在多头注意力结构中，上述过程会在多个子空间里同时进行，然后再把结果拼接起来做线性映射。该机制使模型能够从各个表示子空间中获得多样化的语义关系。多头注意力机制在模型结构里主要处在编码器和解码器层中，该模块对建模全局依赖关系起着主要作用。处理长序列时，时间复杂度和空间复杂度都会达到 $O(L^{2})$，这对显存带宽和片上缓存的要求较高。特别是中间生成的相关性矩阵规模较大，如果频繁读写全局内存，会严重制约整体性能\cite{beltagy2020longformer}。

% 第三节
\section{机器学习编译与 MLIR 基础}

\subsection{机器学习编译器概述}
随着深度学习模型复杂度的不断提升与硬件后端的日益多样化，传统的基于库的算子实现方式逐渐显露出局限性。为了解决模型与硬件之间高效映射的难题，机器学习编译器\cite{zhang2023compiler}应运而生。机器学习编译器是一种专门用于将高层次的深度学习模型描述转换为特定硬件上可执行的高性能机器代码的系统软件。其核心目标是弥合上层算法框架与底层异构硬件之间的语义鸿沟，在减轻人工优化负担的同时，实现代码的自动化生成与性能调优。如图 \ref{2_7} 所示，传统编译器与机器学习编译器在三个组成部分上存在差异，其中一个成熟的机器学习编译器架构通常包含三个主要阶段。

\begin{figure}[htbp]
\centering
\includegraphics[width=0.85\textwidth]{2_7}
\caption{传统编译器（右）和机器学习编译器（左）的编译流程对比}\label{2_7}
\vspace{\baselineskip}
\end{figure}

首先是前端和图级优化，编译器先会把深度学习框架定义的计算图转换成高层中间表示\cite{roesch2018relay, zhong2025flashtensor}。编译器在该阶段主要负责执行和硬件无关的优化操作，具体包括算子融合\cite{JSJA202505006}、代数化简、死代码消除。算子融合把多个计算密集型算子和访存密集型算子合在一起执行，能够明显减少对全局显存的访问开销。

其次是算子级优化和调度，编译器把计算图进行转换，使其变为张量级或者循环级的 IR\cite{feng2023tensorir}。像 TVM 这类典型系统通过建立搜索空间或者使用模板化生成方法，能够自动对循环切分、向量化、并行调度策略进行搜索。这个过程主要用来替代通常的手工编写 CUDA 内核方式，以此适应不同硬件的存储层级和计算特性。

最后是后端与代码生成，优化后的 IR 最终被传递给后端，生成面向特定目标硬件的源代码（如 CUDA、LLVM IR 或 C++），并编译为二进制可执行文件。

虽然现有深度学习编译器在一定程度上完成了多后端支持和自动化优化，但是正如本文前文所述，包括 TVM 在内的较多系统在中间表示层面没有直接刻画 GPU 线程块级结构，难以灵活表达复杂的协同优化策略，给 Transformer 类大模型的高性能算子生成造成了困难。

\subsection{MLIR 核心设计原理}

MLIR 是由 Google 主导建立的新型编译器架构\cite{lattner2021mlir, lattner2020mlir}，用来解决深度学习和异构计算领域内编译器基础设施碎片化日益严重的问题。传统编译模式通常依赖单一抽象层级的中间表示，而 MLIR 提出以可扩展、多层次中间表示为中心的统一编译模式。图 \ref{2_8} 显示了使用 Affine 和 Std 方言进行多项式乘法的 MLIR 通用表示。该图说明 MLIR 的设计思想重点在于让不同抽象层级和不同领域语义的中间表示可以在统一的编译流程中共存并发挥作用。该方法支持从高层算法描述到底层硬件指令的逐步语义转换和优化，在提高优化能力的同时，也降低了跨平台编译系统的工程复杂度。
\begin{figure}[htbp]
\centering
\includegraphics[width=0.72\textwidth]{2_8}
\caption{MLIR 的通用表示\cite{lattner2020mlir}}\label{2_8}
\vspace{\baselineskip}
\end{figure}

MLIR 的方言机制让开发者在特定领域内对操作、类型和属性进行定制，从而形成一种灵活性很高的中间表示系统\cite{data_flow_mlir}。该设计在横向方面实现了多种方言的共存和共同运作，例如用于处理线性代数的 Linalg 方言可以和 GPU 并行计算方言在同一个模块中共同工作，形成了各个领域之间的解耦。MLIR 在纵向维度上提供了多种内置方言，包括从高层图表示 TOSA 到中层循环结构 Affine 和 SCF、到底层硬件抽象 LLVM IR 的完整抽象层级，由此支持渐进式编译优化的结构化需求。多个层级和多个领域构成的方言机制，能让编译器精确描述各阶段的计算语义，并且给复杂的优化工作提供充分的上下文信息。

MLIR 利用逐级降级和转换方式，把高层抽象逐步细化成低层表示\cite{chelini2021progressive}，这种渐进式转换策略和传统编译器一步到位式降级相比，拥有明显的优势。在转换期间，循环边界、张量形状、并行语义等高层语义信息能够得到保留，从而使编译器可以在不同抽象层级上执行针对性优化。例如在循环层级进行循环融合、在张量层级进行算子融合，并且不会因为过早降级而造成语义丢失。该机制也促进优化逻辑进行复用。各种深度学习框架可以共享中间的方言定义和优化 Pass，避免重复开发，提高整个生态的工程效率。

为了提高编译器扩展性、开发效率与可维护性，MLIR 提出了操作定义规范（Operational Definition Specification, ODS）和声明式重写规则（Declarative Rewrite Rules, DRR），并利用 TableGen 技术对方言操作与重写规则进行定义。ODS 允许开发者用声明式方式来定义操作的操作数、结果、属性等元信息，并且会自动生成类型检查、序列化等基础设施代码。DRR 利用模式匹配和重写规则来描述优化转换逻辑，使优化 Pass 的编写变得更直观，也方便进行验证。声明式方法可以促使编译器在扩展过程中维持类型安全和结构一致性，并减少学习新方言和优化规则的成本，这有利于支持 MLIR 生态的发展。

\subsection{基于 MLIR 的代码生成流程}

依靠 MLIR 进行代码生成，实际上是通过多抽象层级中间表示，利用多个结构化的转换和优化 Pass，把高层计算模型逐步降级，最后将其映射到目标硬件。传统编译器利用单一中间表示进行降级，MLIR 的分层设计在编译各个阶段保留最适合的语义信息，优化灵活性和可扩展性因此得到明显提高。图 \ref{2_9} 显示，整个流程整体上可以划分成高层语义建模、中层结构化表示、低层硬件映射、后端代码生成这四个主要阶段。本文把模型优化集成在多层 IR 的降级过程中，而且各级 IR 之间能够混合使用 MLIR 的方言，这显著提高了模型优化的灵活性。

\begin{figure}[htbp]
\centering
\includegraphics[width=0.65\textwidth]{2_9}
\caption{MLIR 编译流程}\label{2_9}
\vspace{\baselineskip}
\end{figure}

进行高层语义建模时，输入一般来自深度学习框架的计算图前端或模型脚本。该阶段的主要目标是将算子和数据依赖关系转换为能够表达计算语义的 MLIR 高层方言表示，具体包括张量级操作和算子级数据流。此类表示方式有利于在抽象层面进行算子融合和代数变换等全局优化，并防止过早引入与硬件相关的细节，从而给后续优化保留完整的语义搜索空间。

随后，流程进入中层结构化表示及调度表达阶段。高层算子语义被扩展为更结构化的循环、分块和并行映射形式，用来解决计算组织和调度方式的问题。该阶段通常通过方言转换和定制化 Pass，把张量计算映射为多级循环结构，并且建立计算任务和线程层次中的抽象关系。中间表示在保留结构化语义的同时，也对存储层级、并行粒度和访存模式进行初步描述，为后续的硬件映射打下基础。

编译器在面向硬件的降级和资源映射阶段，负责把中层表示里的抽象并行结构、存储模型，映射到具体硬件的执行模型和资源系统上。该过程利用逐步引入层级较低的方言对硬件细节进行显式处理，包括线程映射、地址空间、同步语义和寄存器使用。MLIR 的渐进式降级方式让多个阶段性的 Pass 共同发挥作用，该通用流程可以在不同硬件平台上复用，仅需在关键位置加入带有平台特点的个性化转换。

最终，在后端代码生成阶段，经过充分降级和硬件映射的中间表示被转换为 LLVM 等传统后端所支持的低层表示，并生成可执行的目标代码。MLIR 方面已经完成了大部分结构化优化工作，后端主要承担指令选择和寄存器分配等经典编译任务，以此实现高层算子优化与底层指令生成的责任分离，编译流程的稳定性和可维护性由此得到提高。


% 小结
\section{本章小结}

本章首先说明异构计算加速平台发展现状，主要分析主流 NVIDIA GPU 和国产海光 DCU 的硬件架构和编程模型。
首先这两者虽然都符合 SIMT 执行模型，但在主要调度单位和存储层级特性上表现出明显差异，该差异构成了跨平台算子优化的底层困难。

其次本章分析了 Transformer 模型的架构组成和计算特性。本文分析了矩阵乘法和注意力机制这两个主要算子，发现大模型推理呈现出计算密集和访存密集交织的特点，说明对这些算子进行深度优化是必要的。

最后分析了机器学习编译器的基本原理，并着重说明了 MLIR 的主要设计思想。MLIR 通过运用多层级方言机制和渐进式降级过程，能够缓解异构硬件生态碎片化问题，并且提高算子生成的效率，从而提供必要的底层基础设施。充分把握异构硬件的微架构特点、深入了解 Transformer 算子的数据流向，并且运用 MLIR 分层编译技术，可以形成高性能和跨平台算子优化策略的基础。

本章的研究为后续章节提出针对性编译优化策略和自定义方言设计提供了有效的理论支持。
